\section{注意力机制：从 RNN 瓶颈到通用操作}

\subsection{注意力机制的基本流程}

注意力机制对 Encoder-Decoder 架构的改进如下：

\begin{figure}[H]
\centering
\includegraphics[width=0.95\textwidth]{slides-images/slide-005.jpg}
\caption{引入注意力机制的 Seq2Seq 模型：编码器 RNN 保持不变，但解码器在每个时间步都可以``回头看''编码器的所有隐藏状态\protect\footnotemark}
\end{figure}
\footnotetext{Slides 第6页。}

以解码器生成第一个输出词为例，完整流程：

\textbf{Step 1：计算对齐分数}（Alignment Scores）

对于解码器的隐藏状态 $s_0$ 和编码器的每个隐藏状态 $h_i$，计算它们的``匹配程度''：

$$e_{i} = f_{\text{att}}(s_0, h_i)$$

其中 $f_{\text{att}}$ 是一个可学习的函数（最初使用的是一个线性层：将 $[s_0; h_i]$ 拼接后通过线性变换映射为标量）。

\textbf{Step 2：Softmax 归一化}

$$a_i = \frac{\exp(e_i)}{\sum_j \exp(e_j)}$$

得到一个在所有输入 token 上的\textbf{概率分布}——这就是\textbf{注意力权重}。

\begin{figure}[H]
\centering
\includegraphics[width=0.95\textwidth]{slides-images/slide-008.jpg}
\caption{注意力权重计算：对齐分数经过 softmax 后成为概率分布，每个值表示解码器在当前时间步对输入序列各位置的``关注度''\protect\footnotemark}
\end{figure}
\footnotetext{Slides 第9页。}

\textbf{Step 3：加权求和}

$$c_1 = \sum_i a_i \cdot h_i$$

注意力权重加权的编码器隐藏状态的线性组合，产生\textbf{上下文向量} $c_1$。

\begin{figure}[H]
\centering
\includegraphics[width=0.95\textwidth]{slides-images/slide-009.jpg}
\caption{上下文向量 $c_1$ 是编码器隐藏状态的加权线性组合，权重由注意力分数决定\protect\footnotemark}
\end{figure}
\footnotetext{Slides 第10页。}

\textbf{Step 4：生成输出}

将上下文向量 $c_1$ 与解码器的其他输入一起送入解码器 RNN，生成下一个隐藏状态和输出 token。

\begin{figure}[H]
\centering
\includegraphics[width=0.95\textwidth]{slides-images/slide-010.jpg}
\caption{每个解码器时间步都执行一次完整的注意力计算：计算新的对齐分数 $\rightarrow$ softmax $\rightarrow$ 加权求和 $\rightarrow$ 生成上下文向量 $\rightarrow$ 输入解码器\protect\footnotemark}
\end{figure}
\footnotetext{Slides 第11页。}

\begin{importantbox}{注意力机制的核心思想}
\begin{enumerate}
    \item 不再通过单一向量传递所有信息
    \item 解码器在\textbf{每个时间步}都重新``审视''整个输入序列
    \item 每次生成不同的上下文向量——动态聚焦于输入中\textbf{与当前输出最相关的部分}
    \item 整个过程\textbf{完全可微}，注意力权重通过梯度下降端到端学习
\end{enumerate}
\end{importantbox}

\subsection{注意力权重的可视化}

\begin{figure}[H]
\centering
\includegraphics[width=0.95\textwidth]{slides-images/slide-014.jpg}
\caption{英法翻译任务中的注意力权重矩阵：行对应输出序列（法语），列对应输入序列（英语）。对角线结构表示词对词的对应关系，反对角线表示词序的调整\protect\footnotemark}
\end{figure}
\footnotetext{Slides 第15页。来自 Bahdanau et al., 2015。}

注意力矩阵的可视化揭示了网络学到的输入-输出对齐关系：
\begin{itemize}
    \item \textbf{对角线结构}：输入和输出中词序一致的部分（如``the agreement'' $\leftrightarrow$ ``L'accord''）
    \item \textbf{反对角线}：词序发生变化的部分（如``European economic area'' 在法语中的不同语序）
    \item \textbf{$2 \times 2$ 块}：一个输入词对应多个输出词或反之
\end{itemize}

\begin{knowledgebox}{注意力机制的历史地位}
注意力机制首次出现在 Bahdanau et al. 2015 年的论文``Neural Machine Translation by Jointly Learning to Align and Translate''中。该论文获得了 ICLR 2025 的 runner-up Test of Time Award。这一工作虽然是在 RNN 框架内提出的，但其核心思想远超 RNN 本身——它开创了一种全新的计算范式。
\end{knowledgebox}

\subsection{本章小结}

注意力机制通过让解码器在每个时间步动态地关注输入序列的不同部分，消除了固定长度上下文向量的瓶颈。这一机制完全可微，通过端到端训练自动学习输入-输出之间的对齐关系。

%% ========== 通用注意力算子 ==========

